% Trois spectres infrarouges à attribuer (exercice 7, chapitre 11)
\begin{tikzpicture}
\begin{axis}[otIR,width=100mm,height=31mm,title={spectre A},xlabel={}]
\addplot[otSpectre] {100 \bandeIR{2960}{55}{42} \bandeIR{1715}{35}{72} \bandeIR{1360}{45}{30} \bandeIR{1220}{40}{28}};
\end{axis}
\end{tikzpicture}
\par\vspace{0.5mm}
\begin{tikzpicture}
\begin{axis}[otIR,width=100mm,height=31mm,title={spectre B},xlabel={}]
\addplot[otSpectre] {100 \bandeIR{3340}{170}{60} \bandeIR{2940}{55}{46} \bandeIR{1420}{55}{20} \bandeIR{1060}{45}{52}};
\end{axis}
\end{tikzpicture}
\par\vspace{0.5mm}
\begin{tikzpicture}
\begin{axis}[otIR,width=100mm,height=31mm,title={spectre C}]
\addplot[otSpectre] {100 \bandeIR{2900}{330}{52} \bandeIR{2960}{60}{18} \bandeIR{1710}{35}{70} \bandeIR{1240}{45}{35}};
\end{axis}
\end{tikzpicture}
\par\smallskip\legende{Les spectres infrarouges des trois liquides, dans le désordre.}
